Un fluido tiene velocidad $\mathbf v=\langle-\alpha y,\alpha x,\beta z\rangle$, donde $\alpha$ y $\beta$ son constantes. Determine la circulación sobre la intersección de $x^2+y^2=a^2$, $a>0$, con $z=h+mx+ny$, orientada con proyección antihoraria sobre $xy$. Interprete qué parámetros influyen en la circulación y relacione el resultado con la vorticidad.
